\noindent\textbf{A source-defined year, with two chronological layers.}
This chapter covers all 46 catalogue records 667--712 and the 128 supplied
pages P001257--P001384. It follows manuscript order while retaining catalogue
record numbers and dates; these are identifiers, not a silently corrected chronology. The manuscript
dates in the opening spring sequence often say 1986 where the catalogue
says 1987, and record 667 has an earlier handwritten year. Such conflicts
are recorded at the affected talks. Shared page ranges are not copied
wholesale into both talks: headings, page layout, and the mathematical
continuity determine the attributed passage. In particular, the Gurevich
heading on P001303 qualifies record 684's catalogue attribution, and the
Beilinson heading on P001319 recovers surviving material for record 694,
whose catalogue entry has no linked pages. A boundary fragment at P001257
is linked to its earlier-year occurrence rather than reassigned to
Kontsevich. The catalogue itself was checked
\citep{mccme2026gelfandcatalogue}.

\medskip
\noindent\textbf{The mathematical lines.}
String measures and determinants, arrangements of real and complex
subspaces, quantum groups and their representations, matrix integrals
and map enumeration, integrable systems, modular lattices, automorphic
forms, and cohomological constructions run through the packet. The
apparently separate talks often share a concrete obstruction: deciding
what remains valid after a deformation, identifying the correct quotient
or completion, or distinguishing a formal generating series from an
analytic integral. The later fate is correspondingly heterogeneous.
Several talks can be matched to near-contemporary publications; others
restate earlier results; a few questions remain meaningful only after
their missing hypotheses and intended output are specified.

\medskip
\noindent\textbf{Inventory and assessment.}
There are 349 distinct indexed items, 21
problem-index entries, and 129 corrections or provenance notices.
An item is a mathematical unit, not a page or a sentence.
\begin{center}
\begin{tabular}{cl r}
\hline
Class & Historical assessment & Items\\
\hline
A & Published soon after by the speaker or Gelfand & 39 \\
B & Developed by others & 24 \\
C & Partly developed & 10 \\
D & No identified development after a documented search & 0 \\
E & Known before the talk & 242 \\
F & Not identifiable from the surviving notes & 34 \\
\hline
\end{tabular}
\end{center}
Class records historical fate; it does not certify literal correctness.
The separate fields retain missing hypotheses, false literal formulas,
partial publication matches, and the actual level of source inspection.
F includes unrecoverable fragments and unsupported precise readings;
it does not mean that their surrounding subject is unknown.
No item is assigned D: the searches did not justify an assertion that a precise idea had received no development.

\medskip
\noindent\textbf{Bounded developments.}
Two class C questions support explicit work. For the quantum KdV spectral
problem, a complete level-two calculation exhibits a genuine Jordan block
in an irreducible, nonunitary Virasoro module. Thus an unrestricted
eigenbasis formulation needs an extra hypothesis; the notebook itself
does not assert semisimplicity. For the two-source unitary integral,
probability Haar normalization gives a regular entire \(U(2)\) formula,
an all-rank convergent moment expansion with a proved tail bound, and
an exact \(U(3)\) counterexample to a specified naive permutation-only
extension. These do not settle an all-grade spectrum or a generic
large-\(N\) asymptotic problem. Two further annotations rederive a known
transverse-plane criterion and certify the notebook's existing answer
of 96 cube regions. Independent derivations or algorithms agree in
every reported test; code and complete outputs accompany the chapter.

\medskip
\noindent\textbf{Questions and continuity.}
The problem index separates construction of quantum KdV charges from
their spectral classification; finite multi-source integrals from the
broader reduced gauge integral; and established special cases from
unrestricted arithmetic conjectures. Earlier problem identifiers are
reused only for exact recurrences, notably the determinant question
GP-1982-40. Other earlier occurrences are linked as antecedents where
their quantifiers differ. Broad thematic resemblance is not an identity.
The toroidal reformulation of the Riemann hypothesis requires the
correct spectral condition, including temperedness; neither an
unrestricted unitarity slogan nor a related theorem settles RH.

\medskip
\noindent\textbf{Status and evidence.}
The cutoff is 12 November 2026. Newly circulated 2026 papers are explicitly
identified as preprints, with the version and theorem actually inspected;
their existence is not presented as independent proof certification.
Every bibliography entry carries a verification note, and unresolved
leads remain distinguishable from results read in the primary source.
The report provides counts by speaker and area, records the difficult
attributions and searches, and specifies how a website should expose
source, reconstruction, literature, and present status separately.
